\clearpage
\section{总结}

本项目完成了一款面向 FPGA 的五级顺序单发射 RISC-V 处理器。处理器以 ready/valid 反压和 epoch
重定向维持顺序执行，通过 512 项 8-bank BTB、2-bit 动态方向预测和返回地址栈改善控制流，通过
两项固定槽取指缓冲隔离译码反压，通过总容量 4\,KiB 的双 bank D-cache、同步普通加载通路和链表查询
镜像副本服务两类不同访存需求，并以分区结果通路控制后端选择器深度和扇出。CSR 读改写局部化到
EXU，两阶段事务在执行级完成读取、准备和顺序提交。

指令能力覆盖 RV32I、RV32M、Zicsr、机器态控制流以及部分 RISC-V B 扩展。自定义指令包括 CRC 字节
更新、融合矩阵乘法/乘累加、链式查询与反转、矩阵归约和流式状态处理；各单元均纳入处理器的握手、相关检测、
访存一致性和顺序写回体系。GCC 工具链能够生成相应自定义指令，使硬件能力以统一的指令接口使用。

总决赛保留链的真实板测时间由 10.695831\,s 降至 5.346451\,s，缩短约 50.01\%。最终合入提交
\code{8591368} 完成流水边界、缓存查询、自定义加速和 GCC 中端/后端生成路径的整合；其 RTL 在
240\,MHz 下 routed WNS 为 $+0.014$\,ns 并完成时序收敛，在 300\,MHz 下 routed WNS 为
$-0.699$\,ns。最终程序的 10/100 次 NPC 仿射拟合为 5.292462350\,s，
NEMU 10,000 次回归 CRC 正确，但该时间尚未上板，文中始终以青色与真实板测分开。

验证体系覆盖 Vivado Simulation、Verilator/NEMU 差分测试、指令与体系结构测试、CoreMark、
Dhrystone、RT-Thread Nano 3.1.5、Vivado 实现和真实 FPGA 板测。由程序镜像、实现报告、bitstream 身份、UART
原始输出和 CRC 组成的证据链，使功能、性能与物理实现结果能够在统一口径下复核。对只改变编译器和
程序镜像、未改变 RTL 的最终迁移，复用已有实现 WNS 并明确标注未重跑 Vivado；没有板级 capture 的
拟合结果不表述为板测结论。
